% Сборка: lualatex lab3-report.tex (рисунки — node ../raster-lab/scripts/report-figures.cjs)
\documentclass[11pt,a4paper]{article}
\usepackage[a4paper,margin=2cm,top=1.6cm,bottom=1.8cm]{geometry}
\usepackage{fontspec}
\usepackage[russian]{babel}
\usepackage{amsmath}
\usepackage{unicode-math}
\setmainfont{Liberation Serif}
\setmonofont{Liberation Mono}[Scale=MatchLowercase]
\setmathfont{TeX Gyre Termes Math}
\usepackage{titlesec}
\usepackage{enumitem}
\usepackage{microtype}
\usepackage{booktabs}
\usepackage{caption}
\usepackage{float}
\usepackage{tikz}
\usepackage{subcaption}

\titleformat{\section}{\large\bfseries}{\thesection.}{0.5em}{}
\titleformat{\subsection}{\normalsize\bfseries}{\thesubsection.}{0.5em}{}
\titlespacing*{\section}{0pt}{1.0ex plus .3ex}{0.5ex}
\titlespacing*{\subsection}{0pt}{0.8ex plus .3ex}{0.4ex}
\setlength{\parindent}{0pt}
\setlength{\parskip}{0.45em}
\pagestyle{plain}
\captionsetup{font=small,labelsep=period,skip=4pt}

\definecolor{pix}{HTML}{6A5CE0}
\definecolor{gridc}{HTML}{B8BDCC}
\definecolor{idealc}{HTML}{E5484D}

\usepackage[hidelinks,pdftitle={Лабораторная работа № 3. Базовые растровые алгоритмы},pdfauthor={Aliaksei Charnukha}]{hyperref}

\begin{document}

\begin{center}
  {\LARGE\bfseries Лабораторная работа № 3}\\[0.4em]
  {\large Базовые растровые алгоритмы}
\end{center}

\section{Цель}
Реализовать и сравнить алгоритмы растеризации: пошаговый, ЦДА, алгоритм
Брезенхема для отрезка и для окружности; дополнительно~— алгоритм
Кастла~— Питвея и сглаженную линию (алгоритм Ву). Показать работу каждого на
примере с промежуточными вычислениями и измерить время работы.

\section{Теория}

\subsection{Растровая сетка и привязка координат}
Пиксель $(x, y)$ с целыми $x, y$~— квадрат $[x - \tfrac12, x + \tfrac12] \times [y - \tfrac12, y + \tfrac12]$,
то есть \textbf{целая точка~— центр клетки}, а линии сетки проходят по полуцелым
координатам. Оси проходят через центры клеток столбца $x = 0$ и строки $y = 0$, ось
$y$ направлена вверх. Растеризовать фигуру~— выбрать клетки, центры которых
ближе всего к ней (рис.~\ref{fig:examples}).

Отрезок $(x_0, y_0)$--$(x_1, y_1)$, $\Delta x = x_1 - x_0$, $\Delta y = y_1 - y_0$.
Если $|\Delta x| \ge |\Delta y|$ (отрезок «пологий»), то на каждый $x$ нужен ровно один
пиксель: шагаем по $x$ с шагом 1 и подбираем $y$. Иначе оси меняются ролями.
Пикселей всегда $\max(|\Delta x|, |\Delta y|) + 1$. Формулы ниже~— для
пологого отрезка с $\Delta x > 0$.

Сквозной пример для отрезков: $(1, 1)$--$(8, 4)$, $\Delta x = 7$, $\Delta y = 3$,
$k = 3/7 \approx 0.4286$.

\subsection{Пошаговый алгоритм}
Прямая $y = kx + b$, $k = \Delta y / \Delta x$, $b = y_0 - k x_0$. Для каждого
$x = x_0, \dots, x_1$ вычисляется $y$ и округляется до ближайшего целого. На каждом
шаге~— умножение и вещественные числа. В примере $k = 3/7$, $b = 1 - 3/7 = 4/7$.

\subsection{ЦДА (цифровой дифференциальный анализатор)}
$L = \max(|\Delta x|, |\Delta y|)$ шагов, приращения $dx = \Delta x / L$, $dy = \Delta y / L$
(одно из них равно $\pm1$). Начиная с $(x_0, y_0)$, на каждом шаге $x \mathrel{+}= dx$,
$y \mathrel{+}= dy$, пиксель~— округлённая точка. Умножение заменено сложением,
но арифметика вещественная. В примере
$L = 7$, $dx = 1$, $dy = 0.4286$.

У обоих алгоритмов одна таблица, различается только способ получения $y$:

\begin{table}[H]
\centering\small
\caption{Пошаговый алгоритм и ЦДА, отрезок $(1,1)$--$(8,4)$.}
\begin{tabular}{lcccccccc}
\toprule
$x$ & 1 & 2 & 3 & 4 & 5 & 6 & 7 & 8 \\
$y = kx + b$ & 1 & 1.4286 & 1.8571 & 2.2857 & 2.7143 & 3.1429 & 3.5714 & 4 \\
пиксель $y$ & 1 & 1 & 2 & 2 & 3 & 3 & 4 & 4 \\
\bottomrule
\end{tabular}
\end{table}

\subsection{Алгоритм Брезенхема для отрезка}
На каждом шаге $x$ растёт на 1, а $y$ либо остаётся, либо растёт на 1. Пусть $d$~—
высота идеальной прямой над центром последнего пикселя. Сначала $d = 0$; при шаге
по $x$ прямая поднимается на $k$: $d \mathrel{+}= k$. Если $d > \tfrac12$, прямая ближе
к верхней клетке: $y \mathrel{+}= 1$ и $d \mathrel{-}= 1$ (высота теперь отсчитывается
от нового центра).

Важен только знак $d - \tfrac12$; умножив на $2\Delta x$, получим целую
\textbf{ошибку} $e = 2\Delta x\,(d - \tfrac12)$:
\[
  e_0 = 2\Delta y - \Delta x,\qquad
  e \ge 0:\ y \mathrel{+}= 1,\ e \mathrel{-}= 2\Delta x;\qquad
  \text{на каждом шаге } e \mathrel{+}= 2\Delta y.
\]
В примере $2\Delta y = 6$,
$2\Delta x = 14$, $e_0 = -1$, $e = 14d - 7$:

\begin{table}[H]
\centering\small
\caption{Алгоритм Брезенхема, отрезок $(1,1)$--$(8,4)$.}
\begin{tabular}{cccccc}
\toprule
Шаг & Пиксель & $d$ & $e$ & Решение & $e$ на следующем шаге \\
\midrule
0 & $(1,1)$ & $3/7$ & $-1$ & $e < 0$: только $x$ & $-1 + 6 = 5$ \\
1 & $(2,1)$ & $6/7$ & $5$ & $e \ge 0$: $x$ и $y$ & $5 - 14 + 6 = -3$ \\
2 & $(3,2)$ & $2/7$ & $-3$ & только $x$ & $3$ \\
3 & $(4,2)$ & $5/7$ & $3$ & $x$ и $y$ & $-5$ \\
4 & $(5,3)$ & $1/7$ & $-5$ & только $x$ & $1$ \\
5 & $(6,3)$ & $4/7$ & $1$ & $x$ и $y$ & $-7$ \\
6 & $(7,4)$ & $0$ & $-7$ & только $x$ & $-1$ \\
7 & $(8,4)$ & & & конец & \\
\bottomrule
\end{tabular}
\end{table}
$d$~— высота прямой над текущим пикселем после следующего шага по $x$.

\subsection{Алгоритм Брезенхема для окружности}
Строится четверть окружности радиуса $R$ с центром в начале координат от $(0, R)$ до
$(R, 0)$; остальные три четверти~— отражением $(\pm x, \pm y)$, затем сдвиг на центр
$(x_c, y_c)$. Из пикселя $(x, y)$ возможны три хода: \textbf{H}~— $(x+1, y)$,
\textbf{D}~— $(x+1, y-1)$, \textbf{V}~— $(x, y-1)$. Ошибка пикселя~— $\varepsilon =
x^2 + y^2 - R^2$ ($>0$~— центр снаружи, $<0$~— внутри); выбирается ход с меньшей
$|\varepsilon|$.

$\Delta = (x+1)^2 + (y-1)^2 - R^2$~— ошибка диагонального пикселя, $\Delta_0 = 2(1 - R)$.
\begin{itemize}[label=--, leftmargin=1.5em, itemsep=0.1em, topsep=0.2em]
  \item $\Delta < 0$ (D внутри, окружность проходит правее и выше): выбор между H и D.
        $\delta = \varepsilon_H + \varepsilon_D = 2\Delta + 2y - 1$; $\delta \le 0 \Rightarrow$ H, иначе D.
  \item $\Delta > 0$ (D снаружи): выбор между D и V. $\delta' = \varepsilon_D + \varepsilon_V =
        2\Delta - 2x - 1$; $\delta' \le 0 \Rightarrow$ D, иначе V.
  \item $\Delta = 0$: D лежит на окружности.
\end{itemize}
Знак суммы ошибок разных знаков показывает, какой кандидат ближе. Пересчёт
без возведения в квадрат (новые $x, y$): H: $\Delta \mathrel{+}= 2x + 1$; D:
$\Delta \mathrel{+}= 2x - 2y + 2$; V: $\Delta \mathrel{+}= -2y + 1$.

\begin{table}[H]
\centering\small
\caption{Алгоритм Брезенхема для окружности, $R = 5$, центр $(0, 0)$.}
\begin{tabular}{cccccl}
\toprule
Шаг & $(x, y)$ & $\Delta$ & $\delta$ / $\delta'$ & Ход & Проверка: $\varepsilon$ кандидатов \\
\midrule
0 & $(0,5)$ & $-8$ & $\delta = -7$ & H & $\varepsilon_H(1,5) = 1$, $\varepsilon_D(1,4) = -8$ \\
1 & $(1,5)$ & $-5$ & $\delta = -1$ & H & $\varepsilon_H(2,5) = 4$, $\varepsilon_D(2,4) = -5$ \\
2 & $(2,5)$ & $0$ & --- & D & $\varepsilon_D(3,4) = 0$ \\
3 & $(3,4)$ & $0$ & --- & D & $\varepsilon_D(4,3) = 0$ \\
4 & $(4,3)$ & $4$ & $\delta' = -1$ & D & $\varepsilon_D(5,2) = 4$, $\varepsilon_V(4,2) = -5$ \\
5 & $(5,2)$ & $12$ & $\delta' = 13$ & V & $\varepsilon_D(6,1) = 12$, $\varepsilon_V(5,1) = 1$ \\
6 & $(5,1)$ & $11$ & $\delta' = 11$ & V & $\varepsilon_D(6,0) = 11$, $\varepsilon_V(5,0) = 0$ \\
7 & $(5,0)$ & & & конец & \\
\bottomrule
\end{tabular}
\end{table}
Вся окружность~— 28 пикселей (рис.~\ref{fig:examples}).

\subsection{Алгоритм Кастла~— Питвея (дополнительно)}
Пологий отрезок~— слово из $a = \Delta x$ ходов: $s$~— шаг по главной оси, $d$~—
диагональный шаг; ходов $d$ ровно $b = \Delta y$. Алгоритм строит слово с равномерно
распределёнными $d$, склеивая блоки как в алгоритме Евклида:
\begin{center}\small
$x = a - b$, $y = b$, $m_1 = s$, $m_2 = d$;\quad
пока $x \ne y$: если $x > y$, то $x \mathrel{-}= y$, $m_2 = m_1 + m_2^R$, иначе $y \mathrel{-}= x$, $m_1 = m_2 + m_1^R$;\\
ответ~— $(m_2 + m_1^R)$, повторённое $x$ раз ($^R$~— строка задом наперёд).
\end{center}
В слове $x$ блоков $m_1$ и $y$ блоков $m_2$; блоки редкого вида приклеиваются
к блокам частого.

\begin{table}[H]
\centering\small
\caption{Алгоритм Кастла~— Питвея, $a = 7$, $b = 3$.}
\begin{tabular}{cccll}
\toprule
$x$ & $y$ & $m_1$ & $m_2$ & Действие \\
\midrule
4 & 3 & s & d & начало: 4 буквы s, 3 буквы d \\
1 & 3 & s & sd & $x > y$: к каждому d приклеен s \\
1 & 2 & sds & sd & $x < y$ \\
1 & 1 & sdsds & sd & $x < y$ \\
\midrule
\multicolumn{5}{l}{Ответ: $m_2 + m_1^R$ = sd + sdsds = \textbf{sdsdsds}} \\
\bottomrule
\end{tabular}
\end{table}
Ходы от $(1,1)$ дают те же пиксели, что алгоритм Брезенхема.

\subsection{Сглаженная линия: алгоритм Ву (дополнительно)}
Брезенхем закрашивает пиксель целиком или никак, поэтому наклонная линия выглядит
«лесенкой» (алиасинг). У Ву на каждом шаге по главной оси закрашиваются \textbf{два}
соседних пикселя: $\lfloor y \rfloor$ с яркостью $1 - \{y\}$ и $\lfloor y \rfloor + 1$ с
яркостью $\{y\}$, где $\{y\}$~— дробная часть точного $y$. Сумма яркостей в столбце
равна 1, центр тяжести столбца лежит на прямой. В примере при $x = 2$: $y = 1.4286$, пиксель $(2,1)$~— яркость $0.57$,
$(2,2)$~— $0.43$ (рис.~\ref{fig:examples}, справа).

\begin{figure}[H]
\centering
\begin{subfigure}[b]{0.31\linewidth}\centering\resizebox{\linewidth}{!}{\begin{tikzpicture}[x=0.5cm,y=0.5cm]
\fill[pix] (0.5,0.5) rectangle +(1,1);
\fill[pix] (1.5,0.5) rectangle +(1,1);
\fill[pix] (2.5,1.5) rectangle +(1,1);
\fill[pix] (3.5,1.5) rectangle +(1,1);
\fill[pix] (4.5,2.5) rectangle +(1,1);
\fill[pix] (5.5,2.5) rectangle +(1,1);
\fill[pix] (6.5,3.5) rectangle +(1,1);
\fill[pix] (7.5,3.5) rectangle +(1,1);
\draw[gridc,very thin] (-0.5,-0.5) grid[shift={(0.5,0.5)}] (9.5,5.5);
\fill[gray] (0,0) circle[radius=0.05];
\fill[gray] (0,1) circle[radius=0.05];
\fill[gray] (0,2) circle[radius=0.05];
\fill[gray] (0,3) circle[radius=0.05];
\fill[gray] (0,4) circle[radius=0.05];
\fill[gray] (0,5) circle[radius=0.05];
\fill[gray] (1,0) circle[radius=0.05];
\fill[gray] (1,1) circle[radius=0.05];
\fill[gray] (1,2) circle[radius=0.05];
\fill[gray] (1,3) circle[radius=0.05];
\fill[gray] (1,4) circle[radius=0.05];
\fill[gray] (1,5) circle[radius=0.05];
\fill[gray] (2,0) circle[radius=0.05];
\fill[gray] (2,1) circle[radius=0.05];
\fill[gray] (2,2) circle[radius=0.05];
\fill[gray] (2,3) circle[radius=0.05];
\fill[gray] (2,4) circle[radius=0.05];
\fill[gray] (2,5) circle[radius=0.05];
\fill[gray] (3,0) circle[radius=0.05];
\fill[gray] (3,1) circle[radius=0.05];
\fill[gray] (3,2) circle[radius=0.05];
\fill[gray] (3,3) circle[radius=0.05];
\fill[gray] (3,4) circle[radius=0.05];
\fill[gray] (3,5) circle[radius=0.05];
\fill[gray] (4,0) circle[radius=0.05];
\fill[gray] (4,1) circle[radius=0.05];
\fill[gray] (4,2) circle[radius=0.05];
\fill[gray] (4,3) circle[radius=0.05];
\fill[gray] (4,4) circle[radius=0.05];
\fill[gray] (4,5) circle[radius=0.05];
\fill[gray] (5,0) circle[radius=0.05];
\fill[gray] (5,1) circle[radius=0.05];
\fill[gray] (5,2) circle[radius=0.05];
\fill[gray] (5,3) circle[radius=0.05];
\fill[gray] (5,4) circle[radius=0.05];
\fill[gray] (5,5) circle[radius=0.05];
\fill[gray] (6,0) circle[radius=0.05];
\fill[gray] (6,1) circle[radius=0.05];
\fill[gray] (6,2) circle[radius=0.05];
\fill[gray] (6,3) circle[radius=0.05];
\fill[gray] (6,4) circle[radius=0.05];
\fill[gray] (6,5) circle[radius=0.05];
\fill[gray] (7,0) circle[radius=0.05];
\fill[gray] (7,1) circle[radius=0.05];
\fill[gray] (7,2) circle[radius=0.05];
\fill[gray] (7,3) circle[radius=0.05];
\fill[gray] (7,4) circle[radius=0.05];
\fill[gray] (7,5) circle[radius=0.05];
\fill[gray] (8,0) circle[radius=0.05];
\fill[gray] (8,1) circle[radius=0.05];
\fill[gray] (8,2) circle[radius=0.05];
\fill[gray] (8,3) circle[radius=0.05];
\fill[gray] (8,4) circle[radius=0.05];
\fill[gray] (8,5) circle[radius=0.05];
\fill[gray] (9,0) circle[radius=0.05];
\fill[gray] (9,1) circle[radius=0.05];
\fill[gray] (9,2) circle[radius=0.05];
\fill[gray] (9,3) circle[radius=0.05];
\fill[gray] (9,4) circle[radius=0.05];
\fill[gray] (9,5) circle[radius=0.05];
\draw[->,thick] (-0.5,0) -- (9.9,0) node[right] {$x$};
\draw[->,thick] (0,-0.5) -- (0,5.9) node[above] {$y$};
\node[font=\scriptsize,below,fill=white,inner sep=1pt] at (1,-0.15) {$1$};
\node[font=\scriptsize,below,fill=white,inner sep=1pt] at (2,-0.15) {$2$};
\node[font=\scriptsize,below,fill=white,inner sep=1pt] at (3,-0.15) {$3$};
\node[font=\scriptsize,below,fill=white,inner sep=1pt] at (4,-0.15) {$4$};
\node[font=\scriptsize,below,fill=white,inner sep=1pt] at (5,-0.15) {$5$};
\node[font=\scriptsize,below,fill=white,inner sep=1pt] at (6,-0.15) {$6$};
\node[font=\scriptsize,below,fill=white,inner sep=1pt] at (7,-0.15) {$7$};
\node[font=\scriptsize,below,fill=white,inner sep=1pt] at (8,-0.15) {$8$};
\node[font=\scriptsize,below,fill=white,inner sep=1pt] at (9,-0.15) {$9$};
\node[font=\scriptsize,left,fill=white,inner sep=1pt] at (-0.15,1) {$1$};
\node[font=\scriptsize,left,fill=white,inner sep=1pt] at (-0.15,2) {$2$};
\node[font=\scriptsize,left,fill=white,inner sep=1pt] at (-0.15,3) {$3$};
\node[font=\scriptsize,left,fill=white,inner sep=1pt] at (-0.15,4) {$4$};
\node[font=\scriptsize,left,fill=white,inner sep=1pt] at (-0.15,5) {$5$};
\node[font=\scriptsize,below left] at (0,0) {$0$};
\draw[idealc,thick] (1,1) -- (8,4);
\fill[idealc] (1,1) circle[radius=0.15];
\fill[idealc] (8,4) circle[radius=0.15];
\end{tikzpicture}
}\caption{Брезенхем}\end{subfigure}\hfill
\begin{subfigure}[b]{0.31\linewidth}\centering\resizebox{0.85\linewidth}{!}{\begin{tikzpicture}[x=0.42cm,y=0.42cm]
\fill[pix] (-0.5,4.5) rectangle +(1,1);
\fill[pix] (-0.5,-5.5) rectangle +(1,1);
\fill[pix] (0.5,4.5) rectangle +(1,1);
\fill[pix] (-1.5,4.5) rectangle +(1,1);
\fill[pix] (0.5,-5.5) rectangle +(1,1);
\fill[pix] (-1.5,-5.5) rectangle +(1,1);
\fill[pix] (1.5,4.5) rectangle +(1,1);
\fill[pix] (-2.5,4.5) rectangle +(1,1);
\fill[pix] (1.5,-5.5) rectangle +(1,1);
\fill[pix] (-2.5,-5.5) rectangle +(1,1);
\fill[pix] (2.5,3.5) rectangle +(1,1);
\fill[pix] (-3.5,3.5) rectangle +(1,1);
\fill[pix] (2.5,-4.5) rectangle +(1,1);
\fill[pix] (-3.5,-4.5) rectangle +(1,1);
\fill[pix] (3.5,2.5) rectangle +(1,1);
\fill[pix] (-4.5,2.5) rectangle +(1,1);
\fill[pix] (3.5,-3.5) rectangle +(1,1);
\fill[pix] (-4.5,-3.5) rectangle +(1,1);
\fill[pix] (4.5,1.5) rectangle +(1,1);
\fill[pix] (-5.5,1.5) rectangle +(1,1);
\fill[pix] (4.5,-2.5) rectangle +(1,1);
\fill[pix] (-5.5,-2.5) rectangle +(1,1);
\fill[pix] (4.5,0.5) rectangle +(1,1);
\fill[pix] (-5.5,0.5) rectangle +(1,1);
\fill[pix] (4.5,-1.5) rectangle +(1,1);
\fill[pix] (-5.5,-1.5) rectangle +(1,1);
\fill[pix] (4.5,-0.5) rectangle +(1,1);
\fill[pix] (-5.5,-0.5) rectangle +(1,1);
\draw[gridc,very thin] (-6.5,-6.5) grid[shift={(0.5,0.5)}] (6.5,6.5);
\fill[gray] (-6,-6) circle[radius=0.05];
\fill[gray] (-6,-5) circle[radius=0.05];
\fill[gray] (-6,-4) circle[radius=0.05];
\fill[gray] (-6,-3) circle[radius=0.05];
\fill[gray] (-6,-2) circle[radius=0.05];
\fill[gray] (-6,-1) circle[radius=0.05];
\fill[gray] (-6,0) circle[radius=0.05];
\fill[gray] (-6,1) circle[radius=0.05];
\fill[gray] (-6,2) circle[radius=0.05];
\fill[gray] (-6,3) circle[radius=0.05];
\fill[gray] (-6,4) circle[radius=0.05];
\fill[gray] (-6,5) circle[radius=0.05];
\fill[gray] (-6,6) circle[radius=0.05];
\fill[gray] (-5,-6) circle[radius=0.05];
\fill[gray] (-5,-5) circle[radius=0.05];
\fill[gray] (-5,-4) circle[radius=0.05];
\fill[gray] (-5,-3) circle[radius=0.05];
\fill[gray] (-5,-2) circle[radius=0.05];
\fill[gray] (-5,-1) circle[radius=0.05];
\fill[gray] (-5,0) circle[radius=0.05];
\fill[gray] (-5,1) circle[radius=0.05];
\fill[gray] (-5,2) circle[radius=0.05];
\fill[gray] (-5,3) circle[radius=0.05];
\fill[gray] (-5,4) circle[radius=0.05];
\fill[gray] (-5,5) circle[radius=0.05];
\fill[gray] (-5,6) circle[radius=0.05];
\fill[gray] (-4,-6) circle[radius=0.05];
\fill[gray] (-4,-5) circle[radius=0.05];
\fill[gray] (-4,-4) circle[radius=0.05];
\fill[gray] (-4,-3) circle[radius=0.05];
\fill[gray] (-4,-2) circle[radius=0.05];
\fill[gray] (-4,-1) circle[radius=0.05];
\fill[gray] (-4,0) circle[radius=0.05];
\fill[gray] (-4,1) circle[radius=0.05];
\fill[gray] (-4,2) circle[radius=0.05];
\fill[gray] (-4,3) circle[radius=0.05];
\fill[gray] (-4,4) circle[radius=0.05];
\fill[gray] (-4,5) circle[radius=0.05];
\fill[gray] (-4,6) circle[radius=0.05];
\fill[gray] (-3,-6) circle[radius=0.05];
\fill[gray] (-3,-5) circle[radius=0.05];
\fill[gray] (-3,-4) circle[radius=0.05];
\fill[gray] (-3,-3) circle[radius=0.05];
\fill[gray] (-3,-2) circle[radius=0.05];
\fill[gray] (-3,-1) circle[radius=0.05];
\fill[gray] (-3,0) circle[radius=0.05];
\fill[gray] (-3,1) circle[radius=0.05];
\fill[gray] (-3,2) circle[radius=0.05];
\fill[gray] (-3,3) circle[radius=0.05];
\fill[gray] (-3,4) circle[radius=0.05];
\fill[gray] (-3,5) circle[radius=0.05];
\fill[gray] (-3,6) circle[radius=0.05];
\fill[gray] (-2,-6) circle[radius=0.05];
\fill[gray] (-2,-5) circle[radius=0.05];
\fill[gray] (-2,-4) circle[radius=0.05];
\fill[gray] (-2,-3) circle[radius=0.05];
\fill[gray] (-2,-2) circle[radius=0.05];
\fill[gray] (-2,-1) circle[radius=0.05];
\fill[gray] (-2,0) circle[radius=0.05];
\fill[gray] (-2,1) circle[radius=0.05];
\fill[gray] (-2,2) circle[radius=0.05];
\fill[gray] (-2,3) circle[radius=0.05];
\fill[gray] (-2,4) circle[radius=0.05];
\fill[gray] (-2,5) circle[radius=0.05];
\fill[gray] (-2,6) circle[radius=0.05];
\fill[gray] (-1,-6) circle[radius=0.05];
\fill[gray] (-1,-5) circle[radius=0.05];
\fill[gray] (-1,-4) circle[radius=0.05];
\fill[gray] (-1,-3) circle[radius=0.05];
\fill[gray] (-1,-2) circle[radius=0.05];
\fill[gray] (-1,-1) circle[radius=0.05];
\fill[gray] (-1,0) circle[radius=0.05];
\fill[gray] (-1,1) circle[radius=0.05];
\fill[gray] (-1,2) circle[radius=0.05];
\fill[gray] (-1,3) circle[radius=0.05];
\fill[gray] (-1,4) circle[radius=0.05];
\fill[gray] (-1,5) circle[radius=0.05];
\fill[gray] (-1,6) circle[radius=0.05];
\fill[gray] (0,-6) circle[radius=0.05];
\fill[gray] (0,-5) circle[radius=0.05];
\fill[gray] (0,-4) circle[radius=0.05];
\fill[gray] (0,-3) circle[radius=0.05];
\fill[gray] (0,-2) circle[radius=0.05];
\fill[gray] (0,-1) circle[radius=0.05];
\fill[gray] (0,0) circle[radius=0.05];
\fill[gray] (0,1) circle[radius=0.05];
\fill[gray] (0,2) circle[radius=0.05];
\fill[gray] (0,3) circle[radius=0.05];
\fill[gray] (0,4) circle[radius=0.05];
\fill[gray] (0,5) circle[radius=0.05];
\fill[gray] (0,6) circle[radius=0.05];
\fill[gray] (1,-6) circle[radius=0.05];
\fill[gray] (1,-5) circle[radius=0.05];
\fill[gray] (1,-4) circle[radius=0.05];
\fill[gray] (1,-3) circle[radius=0.05];
\fill[gray] (1,-2) circle[radius=0.05];
\fill[gray] (1,-1) circle[radius=0.05];
\fill[gray] (1,0) circle[radius=0.05];
\fill[gray] (1,1) circle[radius=0.05];
\fill[gray] (1,2) circle[radius=0.05];
\fill[gray] (1,3) circle[radius=0.05];
\fill[gray] (1,4) circle[radius=0.05];
\fill[gray] (1,5) circle[radius=0.05];
\fill[gray] (1,6) circle[radius=0.05];
\fill[gray] (2,-6) circle[radius=0.05];
\fill[gray] (2,-5) circle[radius=0.05];
\fill[gray] (2,-4) circle[radius=0.05];
\fill[gray] (2,-3) circle[radius=0.05];
\fill[gray] (2,-2) circle[radius=0.05];
\fill[gray] (2,-1) circle[radius=0.05];
\fill[gray] (2,0) circle[radius=0.05];
\fill[gray] (2,1) circle[radius=0.05];
\fill[gray] (2,2) circle[radius=0.05];
\fill[gray] (2,3) circle[radius=0.05];
\fill[gray] (2,4) circle[radius=0.05];
\fill[gray] (2,5) circle[radius=0.05];
\fill[gray] (2,6) circle[radius=0.05];
\fill[gray] (3,-6) circle[radius=0.05];
\fill[gray] (3,-5) circle[radius=0.05];
\fill[gray] (3,-4) circle[radius=0.05];
\fill[gray] (3,-3) circle[radius=0.05];
\fill[gray] (3,-2) circle[radius=0.05];
\fill[gray] (3,-1) circle[radius=0.05];
\fill[gray] (3,0) circle[radius=0.05];
\fill[gray] (3,1) circle[radius=0.05];
\fill[gray] (3,2) circle[radius=0.05];
\fill[gray] (3,3) circle[radius=0.05];
\fill[gray] (3,4) circle[radius=0.05];
\fill[gray] (3,5) circle[radius=0.05];
\fill[gray] (3,6) circle[radius=0.05];
\fill[gray] (4,-6) circle[radius=0.05];
\fill[gray] (4,-5) circle[radius=0.05];
\fill[gray] (4,-4) circle[radius=0.05];
\fill[gray] (4,-3) circle[radius=0.05];
\fill[gray] (4,-2) circle[radius=0.05];
\fill[gray] (4,-1) circle[radius=0.05];
\fill[gray] (4,0) circle[radius=0.05];
\fill[gray] (4,1) circle[radius=0.05];
\fill[gray] (4,2) circle[radius=0.05];
\fill[gray] (4,3) circle[radius=0.05];
\fill[gray] (4,4) circle[radius=0.05];
\fill[gray] (4,5) circle[radius=0.05];
\fill[gray] (4,6) circle[radius=0.05];
\fill[gray] (5,-6) circle[radius=0.05];
\fill[gray] (5,-5) circle[radius=0.05];
\fill[gray] (5,-4) circle[radius=0.05];
\fill[gray] (5,-3) circle[radius=0.05];
\fill[gray] (5,-2) circle[radius=0.05];
\fill[gray] (5,-1) circle[radius=0.05];
\fill[gray] (5,0) circle[radius=0.05];
\fill[gray] (5,1) circle[radius=0.05];
\fill[gray] (5,2) circle[radius=0.05];
\fill[gray] (5,3) circle[radius=0.05];
\fill[gray] (5,4) circle[radius=0.05];
\fill[gray] (5,5) circle[radius=0.05];
\fill[gray] (5,6) circle[radius=0.05];
\fill[gray] (6,-6) circle[radius=0.05];
\fill[gray] (6,-5) circle[radius=0.05];
\fill[gray] (6,-4) circle[radius=0.05];
\fill[gray] (6,-3) circle[radius=0.05];
\fill[gray] (6,-2) circle[radius=0.05];
\fill[gray] (6,-1) circle[radius=0.05];
\fill[gray] (6,0) circle[radius=0.05];
\fill[gray] (6,1) circle[radius=0.05];
\fill[gray] (6,2) circle[radius=0.05];
\fill[gray] (6,3) circle[radius=0.05];
\fill[gray] (6,4) circle[radius=0.05];
\fill[gray] (6,5) circle[radius=0.05];
\fill[gray] (6,6) circle[radius=0.05];
\draw[->,thick] (-6.5,0) -- (6.9,0) node[right] {$x$};
\draw[->,thick] (0,-6.5) -- (0,6.9) node[above] {$y$};
\node[font=\scriptsize,below,fill=white,inner sep=1pt] at (-6,-0.15) {$-6$};
\node[font=\scriptsize,below,fill=white,inner sep=1pt] at (-4,-0.15) {$-4$};
\node[font=\scriptsize,below,fill=white,inner sep=1pt] at (-2,-0.15) {$-2$};
\node[font=\scriptsize,below,fill=white,inner sep=1pt] at (2,-0.15) {$2$};
\node[font=\scriptsize,below,fill=white,inner sep=1pt] at (4,-0.15) {$4$};
\node[font=\scriptsize,below,fill=white,inner sep=1pt] at (6,-0.15) {$6$};
\node[font=\scriptsize,left,fill=white,inner sep=1pt] at (-0.15,-6) {$-6$};
\node[font=\scriptsize,left,fill=white,inner sep=1pt] at (-0.15,-4) {$-4$};
\node[font=\scriptsize,left,fill=white,inner sep=1pt] at (-0.15,-2) {$-2$};
\node[font=\scriptsize,left,fill=white,inner sep=1pt] at (-0.15,2) {$2$};
\node[font=\scriptsize,left,fill=white,inner sep=1pt] at (-0.15,4) {$4$};
\node[font=\scriptsize,left,fill=white,inner sep=1pt] at (-0.15,6) {$6$};
\node[font=\scriptsize,below left] at (0,0) {$0$};
\draw[idealc,thick] (0,0) circle[radius=5];
\fill[idealc] (0,0) circle[radius=0.15];
\end{tikzpicture}
}\caption{Окружность $R = 5$}\end{subfigure}\hfill
\begin{subfigure}[b]{0.31\linewidth}\centering\resizebox{\linewidth}{!}{\begin{tikzpicture}[x=0.5cm,y=0.5cm]
\fill[black!100] (0.5,0.5) rectangle +(1,1);
\node[font=\tiny,text=white] at (1,1) {1.00};
\fill[black!57] (1.5,0.5) rectangle +(1,1);
\node[font=\tiny,text=white] at (2,1) {.57};
\fill[black!43] (1.5,1.5) rectangle +(1,1);
\node[font=\tiny,text=black] at (2,2) {.43};
\fill[black!14] (2.5,0.5) rectangle +(1,1);
\node[font=\tiny,text=black] at (3,1) {.14};
\fill[black!86] (2.5,1.5) rectangle +(1,1);
\node[font=\tiny,text=white] at (3,2) {.86};
\fill[black!71] (3.5,1.5) rectangle +(1,1);
\node[font=\tiny,text=white] at (4,2) {.71};
\fill[black!29] (3.5,2.5) rectangle +(1,1);
\node[font=\tiny,text=black] at (4,3) {.29};
\fill[black!29] (4.5,1.5) rectangle +(1,1);
\node[font=\tiny,text=black] at (5,2) {.29};
\fill[black!71] (4.5,2.5) rectangle +(1,1);
\node[font=\tiny,text=white] at (5,3) {.71};
\fill[black!86] (5.5,2.5) rectangle +(1,1);
\node[font=\tiny,text=white] at (6,3) {.86};
\fill[black!14] (5.5,3.5) rectangle +(1,1);
\node[font=\tiny,text=black] at (6,4) {.14};
\fill[black!43] (6.5,2.5) rectangle +(1,1);
\node[font=\tiny,text=black] at (7,3) {.43};
\fill[black!57] (6.5,3.5) rectangle +(1,1);
\node[font=\tiny,text=white] at (7,4) {.57};
\fill[black!100] (7.5,3.5) rectangle +(1,1);
\node[font=\tiny,text=white] at (8,4) {1.00};
\draw[gridc,very thin] (-0.5,-0.5) grid[shift={(0.5,0.5)}] (9.5,5.5);
\fill[gray] (0,0) circle[radius=0.05];
\fill[gray] (0,1) circle[radius=0.05];
\fill[gray] (0,2) circle[radius=0.05];
\fill[gray] (0,3) circle[radius=0.05];
\fill[gray] (0,4) circle[radius=0.05];
\fill[gray] (0,5) circle[radius=0.05];
\fill[gray] (1,0) circle[radius=0.05];
\fill[gray] (1,1) circle[radius=0.05];
\fill[gray] (1,2) circle[radius=0.05];
\fill[gray] (1,3) circle[radius=0.05];
\fill[gray] (1,4) circle[radius=0.05];
\fill[gray] (1,5) circle[radius=0.05];
\fill[gray] (2,0) circle[radius=0.05];
\fill[gray] (2,1) circle[radius=0.05];
\fill[gray] (2,2) circle[radius=0.05];
\fill[gray] (2,3) circle[radius=0.05];
\fill[gray] (2,4) circle[radius=0.05];
\fill[gray] (2,5) circle[radius=0.05];
\fill[gray] (3,0) circle[radius=0.05];
\fill[gray] (3,1) circle[radius=0.05];
\fill[gray] (3,2) circle[radius=0.05];
\fill[gray] (3,3) circle[radius=0.05];
\fill[gray] (3,4) circle[radius=0.05];
\fill[gray] (3,5) circle[radius=0.05];
\fill[gray] (4,0) circle[radius=0.05];
\fill[gray] (4,1) circle[radius=0.05];
\fill[gray] (4,2) circle[radius=0.05];
\fill[gray] (4,3) circle[radius=0.05];
\fill[gray] (4,4) circle[radius=0.05];
\fill[gray] (4,5) circle[radius=0.05];
\fill[gray] (5,0) circle[radius=0.05];
\fill[gray] (5,1) circle[radius=0.05];
\fill[gray] (5,2) circle[radius=0.05];
\fill[gray] (5,3) circle[radius=0.05];
\fill[gray] (5,4) circle[radius=0.05];
\fill[gray] (5,5) circle[radius=0.05];
\fill[gray] (6,0) circle[radius=0.05];
\fill[gray] (6,1) circle[radius=0.05];
\fill[gray] (6,2) circle[radius=0.05];
\fill[gray] (6,3) circle[radius=0.05];
\fill[gray] (6,4) circle[radius=0.05];
\fill[gray] (6,5) circle[radius=0.05];
\fill[gray] (7,0) circle[radius=0.05];
\fill[gray] (7,1) circle[radius=0.05];
\fill[gray] (7,2) circle[radius=0.05];
\fill[gray] (7,3) circle[radius=0.05];
\fill[gray] (7,4) circle[radius=0.05];
\fill[gray] (7,5) circle[radius=0.05];
\fill[gray] (8,0) circle[radius=0.05];
\fill[gray] (8,1) circle[radius=0.05];
\fill[gray] (8,2) circle[radius=0.05];
\fill[gray] (8,3) circle[radius=0.05];
\fill[gray] (8,4) circle[radius=0.05];
\fill[gray] (8,5) circle[radius=0.05];
\fill[gray] (9,0) circle[radius=0.05];
\fill[gray] (9,1) circle[radius=0.05];
\fill[gray] (9,2) circle[radius=0.05];
\fill[gray] (9,3) circle[radius=0.05];
\fill[gray] (9,4) circle[radius=0.05];
\fill[gray] (9,5) circle[radius=0.05];
\draw[->,thick] (-0.5,0) -- (9.9,0) node[right] {$x$};
\draw[->,thick] (0,-0.5) -- (0,5.9) node[above] {$y$};
\node[font=\scriptsize,below,fill=white,inner sep=1pt] at (1,-0.15) {$1$};
\node[font=\scriptsize,below,fill=white,inner sep=1pt] at (2,-0.15) {$2$};
\node[font=\scriptsize,below,fill=white,inner sep=1pt] at (3,-0.15) {$3$};
\node[font=\scriptsize,below,fill=white,inner sep=1pt] at (4,-0.15) {$4$};
\node[font=\scriptsize,below,fill=white,inner sep=1pt] at (5,-0.15) {$5$};
\node[font=\scriptsize,below,fill=white,inner sep=1pt] at (6,-0.15) {$6$};
\node[font=\scriptsize,below,fill=white,inner sep=1pt] at (7,-0.15) {$7$};
\node[font=\scriptsize,below,fill=white,inner sep=1pt] at (8,-0.15) {$8$};
\node[font=\scriptsize,below,fill=white,inner sep=1pt] at (9,-0.15) {$9$};
\node[font=\scriptsize,left,fill=white,inner sep=1pt] at (-0.15,1) {$1$};
\node[font=\scriptsize,left,fill=white,inner sep=1pt] at (-0.15,2) {$2$};
\node[font=\scriptsize,left,fill=white,inner sep=1pt] at (-0.15,3) {$3$};
\node[font=\scriptsize,left,fill=white,inner sep=1pt] at (-0.15,4) {$4$};
\node[font=\scriptsize,left,fill=white,inner sep=1pt] at (-0.15,5) {$5$};
\node[font=\scriptsize,below left] at (0,0) {$0$};
\draw[idealc,thick] (1,1) -- (8,4);
\end{tikzpicture}
}\caption{Ву, подписана яркость}\end{subfigure}
\caption{Примеры: точки~— центры пикселей, красное~— идеальная фигура.}
\label{fig:examples}
\end{figure}

\section{Ход работы}
Приложение написано на HTML, CSS и JavaScript и открывается в браузере без
установки (\texttt{lab3/index.html}). Алгоритмы вынесены в модуль
\texttt{raster.js}. Возможности:
\begin{itemize}[label=--, leftmargin=1.5em, itemsep=0.1em, topsep=0.2em]
  \item выбор алгоритма, ввод координат в поля или щелчками по сетке;
  \item оси с делениями и подписями, сетка по границам клеток, центры пикселей,
        идеальная фигура поверх растровой;
  \item масштаб 4--80 точек на клетку, сдвиг перетаскиванием;
  \item таблица вычислений по шагам (как таблицы 1--4), номера шагов или
        яркости в клетках;
  \item замер времени.
\end{itemize}
Модуль проверен 14 юнит-тестами: связность, концы и отклонение от прямой для
312 отрезков, симметрия и минимальность $|x^2 + y^2 - R^2|$ для окружностей
$R = 0 \dots 60$, яркости Ву, ручные примеры из таблиц.

\section{Эксперименты}
Набор~— 20\,000 случайных фигур: концы отрезков и центры окружностей в
квадрате $\pm M$, радиусы 1--$M$. Средняя длина отрезка~— 11, 94 и 931 пиксель.

\subsection{Совпадение результатов}
В «ничьей» прямая проходит ровно посередине между двумя пикселями. Пошаговый
алгоритм и ЦДА округляют половину вверх, Брезенхем~— в сторону движения.

\begin{table}[H]
\centering\small
\caption{Доля отрезков с различающимися пикселями, \%.}
\begin{tabular}{lccc}
\toprule
 & $M = 10$ & $M = 100$ & $M = 1000$ \\
\midrule
Отрезки с «ничьей»       & 33.2 & 32.9 & 33.4 \\
Пошаговый $\ne$ Брезенхем & 16.5 & 16.3 & 16.8 \\
ЦДА $\ne$ Брезенхем       & 16.9 & 17.3 & 17.9 \\
ЦДА $\ne$ пошаговый       & 5.6  & 14.9 & 17.5 \\
Кастл~— Питвей $\ne$ Брезенхем & 0 & 0 & 0 \\
\bottomrule
\end{tabular}
\end{table}
Все расхождения~— на отрезках с «ничьей». Пошаговый и Брезенхем расходятся,
когда вторая координата убывает. Брезенхем от конца к началу даёт другие
пиксели на 33\,\% отрезков. У ЦДА половина получается суммой вещественных
приращений с погрешностью, поэтому с ростом длины он всё чаще расходится с
пошаговым.

\subsection{Время работы}
Пиксели не рисуются, а суммируются; лучший из 15 прогонов после прогрева
JIT (Intel Core i9-14900HX, Node.js 20; разброс 10--20\,\%).

\begin{table}[H]
\centering\small
\caption{Время работы: микросекунды на фигуру / наносекунды на пиксель.}
\begin{tabular}{lccc}
\toprule
Алгоритм & $M = 10$ & $M = 100$ & $M = 1000$ \\
\midrule
Пошаговый                & 0.050 / 4.62 & 0.386 / 4.09 & 3.64 / 3.91 \\
ЦДА                      & 0.034 / 3.17 & 0.258 / 2.73 & 3.99 / 4.29 \\
Брезенхем (отрезок)      & 0.067 / 6.19 & 0.428 / 4.53 & 3.80 / 4.08 \\
Кастл~— Питвей           & 0.512 / 47.7 & 5.63 / 59.7  & 56.8 / 61.0 \\
Ву (сглаживание)         & 0.110 / 6.05 & 0.803 / 4.36 & 7.88 / 4.25 \\
Брезенхем (окружность)   & 0.072 / 2.34 & 0.715 / 2.51 & 7.44 / 2.64 \\
\bottomrule
\end{tabular}
\end{table}
Пошаговый, ЦДА и Брезенхем различаются не более чем в 1.5 раза: в JavaScript
все числа вещественные, и целочисленность Брезенхема не даёт выигрыша.
Окружность дешевле всего на пиксель за счёт симметрии, Ву вдвое медленнее на
фигуру из-за двойного числа пикселей, Кастл~— Питвей в 8--15 раз медленнее из-за
работы со строками.

\section{Результаты}
\begin{itemize}[label=--, leftmargin=1.5em, itemsep=0.2em, topsep=0.2em]
  \item Разработано веб-приложение: пошаговый алгоритм, ЦДА, Брезенхем для
        отрезка и окружности, дополнительно Кастл~— Питвей и Ву; проверено 14 юнит-тестами.
  \item Без «ничьих» все алгоритмы отрезка дают одинаковые пиксели; с ними
        пошаговый и ЦДА расходятся с Брезенхемом на 16--18\,\% отрезков,
        Кастл~— Питвей~— ни на одном.
  \item Время на пиксель: пошаговый, ЦДА, Брезенхем~— 2.7--6.2~нс, Ву~—
        4.3--6.1~нс, окружность~— 2.3--2.6~нс, Кастл~— Питвей~— 48--61~нс.
\end{itemize}

\section{Выводы}
\begin{itemize}[label=--, leftmargin=1.5em, itemsep=0.2em, topsep=0.2em]
  \item Алгоритмы отрезка дают один результат и различаются только
        арифметикой: от пошагового к Брезенхему последовательно убираются
        умножение и вещественные числа.
  \item У трети отрезков есть «ничьи», поэтому правило округления половины
        нужно фиксировать. ЦДА в них ошибается из-за накопленной погрешности,
        Брезенхем точен при любой длине.
  \item Целочисленность Брезенхема ускоряет работу только там, где
        вещественная арифметика дорогая; на современном процессоре скорость
        всех трёх алгоритмов одинакова.
  \item Кастл~— Питвей показывает, что растровая прямая~— периодическое слово
        ходов, строящееся алгоритмом Евклида, но на практике медленнее Брезенхема.
  \item Сглаживание Ву убирает «лесенку» ценой вдвое большего числа пикселей и
        требует оттенков серого.
\end{itemize}

\end{document}
